\section{Preliminaries}\label{sec:preliminaries}

\subsection{Embodied Agent Environments}

We consider an embodied agent with a visual and text environment through a minimal interface: at each timestep $t$, the agent receives a frame observation $o_t$ (a rendered image of the current game state) and selects an action $a_t$ from a fixed set of button inputs $\mathcal{A}$. The environment is partially observable; the agent cannot access fullinternal game state such as NPC positions or battle mechanics beyond what is visually rendered. This formulation is general: any environment that exposes visual output and accepts discrete inputs fits this interface.

\subsection{Agentic Harnesses}

\tersoo{You mentioned RLMs earlier. Currently, we have a trajectory system (initially envisioned for prompt optimization) that basically has the full history of the execution (state, prompt for llm (including reasoning, steps, knowledgebase etc), action,...) perhaps we should also make note of this in our harness... i can edit the summarizer() sub agent to summarize any segment of this window (rather than simply the last n steps (capped at n=50) if you'd like, which would make our harness more akin to RLM. curious about your thought}

\kiran{some of the results in this work are pre-RLMs (I suspect they were inspired by this work significantly as well). they lean into the recursive aspect however, while this approach does not (e.g. it's depth-1).}

An \emph{agentic harness} $\mathcal{H}$ is the scaffolding layer between a foundation model $M$ and the environment. Following the decomposition from~\citet{karten2025pokeagent}, a harness mediates agent behavior through several components:
\begin{itemize}
    \item \textbf{System prompt} $p$: the instructions and strategic guidance provided to the model at each reasoning step.
    \item \textbf{Sub-agents} $\mathcal{S} = \{s_1, \ldots, s_k\}$: specialized modules that can be invoked by the orchestrator for specific tasks (e.g., battle strategy, puzzle solving, self-reflection).
    \item \textbf{Memory} $\mathcal{M}$: a persistent knowledge store that accumulates facts, strategies, and observations across the agent's trajectory.
    \item \textbf{Skills} $\mathcal{K}$: reusable behavioral routines that encode solutions to previously encountered situations.
    \item \textbf{Tools} $\mathcal{T}$: programmatic capabilities available to the model (e.g., pathfinding, information lookup).
\end{itemize}

A \emph{minimalist harness} $\mathcal{H}_{\min}$ provides only the environment interface (frame observations and button inputs) with a generic system prompt and no sub-agents, memory, skills, or domain-specific tools. A \emph{human-designed harness} $\mathcal{H}_{\text{expert}}$ populates all components through manual engineering. A \emph{meta-tool harness} $\mathcal{H}_{\text{meta}}$ takes an intermediate approach: rather than pre-building specific agents, it gives the model meta-tools for creating its own agents, tools, and memory entries (e.g., \texttt{define\_agent}, \texttt{run\_code}, custom tool creation). This was the approach taken in later GPP runs~\citep{zhang2025gpp}, where the model used meta-tools to spontaneously construct pathfinders, battle strategists, and reusable scripts. AutoEvolve starts from $\mathcal{H}_{\min}$ augmented with meta-tools and adds an evolutionary optimization loop that automatically triggers harness refinement based on trajectory analysis, evolving toward $\mathcal{H}_{\text{auto}}$ during a single episode.

\subsection{In-Context Prompt Evolution}

Prompt evolution methods optimize the instructions given to a language model by treating the prompt as a parameter and using trajectory outcomes as signal. GEPA~\citep{agrawal2025gepa} evolves prompts by running complete episodes, scoring the resulting trajectories, and using a meta-model to rewrite the prompt for the next episode. This requires \emph{resetting} the environment between iterations: each evolution step begins from the initial state.

More broadly, Recursive Language Models~\citep{zhang2025rlm} formalize the paradigm of LLMs that programmatically interact with and recurse over their own context, providing a theoretical foundation for agents that self-modify their reasoning process during inference.

AutoEvolve extends this paradigm to the reset-free setting: evolution occurs mid-episode using trajectory \emph{segments} rather than complete episodes, enabling continuous adaptation without interrupting the agent's progress. We detail this mechanism in Section~\ref{sec:methodology}.
